% =============================================================================
% core/paragraph.tex — modelo de parágrafo do Google Docs
% -----------------------------------------------------------------------------
% Todo bloco do currículo é um \cvParagraph. Ele reproduz as regras do Docs:
%
%  * altura de cada linha = line-height × (asc + desc) × tamanho da fonte;
%    a linha de base fica a "asc × tamanho" do topo e a sobra fica embaixo;
%  * espaço entre parágrafos = max(espaço depois do anterior,
%                                  espaço antes do próximo);
%  * quebra de linha gulosa + justificação esticando só os espaços
%    (core/linebreak.lua);
%  * logo após um título de seção, o "espaço antes" do bloco é ignorado.
%
% Uso:
%   \cvParagraph[<chaves>]{<conteúdo>}
%   \cvDefineStyle{<nome>}{<chaves>}   e depois  \cvParagraph[style=<nome>]{...}
%
% Chaves:
%   style          estilo previamente definido (pode ser sobrescrito depois)
%   metrics        sans | serif           métricas verticais da fonte principal
%   font           comando de fonte       ex.: \cvSans, \cvSerifSemiBold
%   size           número (pt do Docs)    ex.: 11
%   line-height    número                 ex.: 1.15
%   before, after  dimensão               espaço antes/depois do parágrafo
%   align          justify | left
%   indent         dimensão               recuo esquerdo do texto
%   right-indent   dimensão               recuo direito do texto
%   marker         conteúdo               marcador (ex.: \cvBullet{8})
%   marker-offset  dimensão               posição do marcador (da margem)
%   color          nome de cor
%   keep-with-next booleano               não separar do próximo bloco
%   heading        booleano               marca o bloco como título de seção
% =============================================================================

\newattribute\cvSnapAttr
\directlua{dofile("core/linebreak.lua")}
\directlua{dofile("core/snap.lua")}

\newdimen\cvLineAbove
\newdimen\cvLineBelow
\newcount\cvGreedy

\ExplSyntaxOn

% --- Métricas ----------------------------------------------------------------
\NewDocumentCommand \cvDefineMetrics { m m m }
  {
    \tl_const:cn { c__cv_asc_#1_tl }  {#2}
    \tl_const:cn { c__cv_desc_#1_tl } {#3}
  }

% --- Chaves ------------------------------------------------------------------
\tl_new:N   \l__cv_metrics_tl
\tl_new:N   \l__cv_font_tl
\tl_new:N   \l__cv_size_tl
\tl_new:N   \l__cv_lh_tl
\dim_new:N  \l__cv_before_dim
\dim_new:N  \l__cv_after_dim
\bool_new:N \l__cv_justify_bool
\dim_new:N  \l__cv_indent_dim
\dim_new:N  \l__cv_right_indent_dim
\tl_new:N   \l__cv_marker_tl
\dim_new:N  \l__cv_marker_offset_dim
\tl_new:N   \l__cv_color_tl
\bool_new:N \l__cv_keep_bool
\bool_new:N \l__cv_heading_bool
\bool_new:N \g__cv_after_heading_bool

\keys_define:nn { cv / par }
  {
    style          .code:n       = \keys_set:nv { cv / par } { l__cv_style_#1_tl } ,
    metrics        .tl_set:N     = \l__cv_metrics_tl ,
    metrics        .initial:n    = sans ,
    font           .tl_set:N     = \l__cv_font_tl ,
    font           .initial:n    = \cvSans ,
    size           .tl_set:N     = \l__cv_size_tl ,
    size           .initial:n    = 11 ,
    line-height    .tl_set:N     = \l__cv_lh_tl ,
    line-height    .initial:n    = 1.15 ,
    before         .dim_set:N    = \l__cv_before_dim ,
    after          .dim_set:N    = \l__cv_after_dim ,
    align          .choice: ,
    align / justify .code:n      = \bool_set_true:N  \l__cv_justify_bool ,
    align / left    .code:n      = \bool_set_false:N \l__cv_justify_bool ,
    align          .initial:n    = left ,
    indent         .dim_set:N    = \l__cv_indent_dim ,
    right-indent   .dim_set:N    = \l__cv_right_indent_dim ,
    marker         .tl_set:N     = \l__cv_marker_tl ,
    marker-offset  .dim_set:N    = \l__cv_marker_offset_dim ,
    color          .tl_set:N     = \l__cv_color_tl ,
    color          .initial:n    = cvText ,
    keep-with-next .bool_set:N   = \l__cv_keep_bool ,
    heading        .bool_set:N   = \l__cv_heading_bool ,
  }

\NewDocumentCommand \cvDefineStyle { m m }
  { \tl_gset:cn { l__cv_style_#1_tl } {#2} }

% --- Parágrafo ---------------------------------------------------------------
\cs_new_protected:Npn \__cv_set_metrics:
  {
    \fp_set:Nn \l_tmpa_fp { \tl_use:c { c__cv_asc_  \l__cv_metrics_tl _tl } }
    \fp_set:Nn \l_tmpb_fp { \tl_use:c { c__cv_desc_ \l__cv_metrics_tl _tl } }
    \dim_set:Nn \cvLineAbove
      { \fp_to_dim:n { \l_tmpa_fp * (\l__cv_size_tl) * 1bp } }
    \dim_set:Nn \cvLineBelow
      {
        \fp_to_dim:n
          { (\l__cv_lh_tl) * (\l_tmpa_fp + \l_tmpb_fp) * (\l__cv_size_tl) * 1bp }
        - \cvLineAbove
      }
  }

\NewDocumentCommand \cvParagraph { O{} +m }
  {
    \par
    \group_begin:
      \keys_set:nn { cv / par } {#1}
      % espaço antes (colapsa com o "depois" do bloco anterior); o primeiro
      % bloco comum depois de um título não usa o próprio "espaço antes"
      \bool_lazy_and:nnF
        { \g__cv_after_heading_bool }
        { ! \l__cv_heading_bool }
        { \addvspace { \l__cv_before_dim } }
      \bool_gset_eq:NN \g__cv_after_heading_bool \l__cv_heading_bool
      \__cv_set_metrics:
      \cvGreedy = 1 \scan_stop:
      \skip_set:Nn \leftskip { \l__cv_indent_dim }
      \bool_if:NTF \l__cv_justify_bool
        { \skip_set:Nn \rightskip { \l__cv_right_indent_dim } }
        { \skip_set:Nn \rightskip { \l__cv_right_indent_dim plus 1fil } }
      \skip_set:Nn \parfillskip { 0pt plus 1fil }
      \l__cv_font_tl
      \cvFontSize { \l__cv_size_tl }
      \noindent
      % a cor é aplicada em modo horizontal: assim nenhum nó de cor vai
      % para a lista vertical (o que impediria o colapso dos espaços)
      \group_begin:
        \color { \l__cv_color_tl }
        \tl_if_empty:NF \l__cv_marker_tl
          {
            \hbox_overlap_left:n
              {
                \hbox_to_wd:nn { \l__cv_indent_dim - \l__cv_marker_offset_dim }
                  { \l__cv_marker_tl \hfil }
              }
          }
        #2
      \group_end:
      \par
      \bool_if:NT \l__cv_keep_bool { \nobreak }
      \addvspace { \l__cv_after_dim }
    \group_end:
  }

% --- Utilitários inline ------------------------------------------------------

% Seleciona o tamanho de fonte como o Docs: o tamanho em pixels (96 dpi) é
% truncado para múltiplos de 1/64 px antes de posicionar os glifos
% (11pt = 14.667px -> 14.65625px). Isso deixa os avanços ~0.07% menores.
\NewDocumentCommand \cvFontSize { m }
  {
    \fontsize
      { \fp_eval:n { floor( (#1) * 256 / 3 ) * 3 / 256 } bp }
      { 0pt }
    \selectfont
  }

% Quebra de linha manual (Shift+Enter no Docs).
\NewDocumentCommand \cvBreak { }
  { \unskip \nobreak \hfil \penalty -10000 \scan_stop: }

% Sublinhado do Docs: traço de 1bp, na cor do texto, ~0.09em abaixo da linha
% de base; extremidades e topo encaixados em pontos inteiros (core/snap.lua).
\box_new:N \l__cv_ul_box
\NewDocumentCommand \cvUnderline { m }
  {
    \mode_leave_vertical:
    \hbox_set:Nn \l__cv_ul_box {#1}
    \hbox_overlap_right:n
      {
        \setattribute \cvSnapAttr { 1 }
        \tex_vrule:D
          width  \box_wd:N \l__cv_ul_box
          height \dim_eval:n { -0.09em + 0.5bp }
          depth  \dim_eval:n {  0.09em + 0.5bp }
      }
    \box_use:N \l__cv_ul_box
  }

% Marcador "●" do Docs (Arial), no tamanho indicado em pt.
\NewDocumentCommand \cvBullet { m O{} }
  {
    \group_begin:
      \tl_if_empty:nF {#2} { \color {#2} }
      \cvBulletFont \cvFontSize {#1} ●
    \group_end:
  }

\ExplSyntaxOff
